\chapter{SQL Foundations with Real OULAD Questions}
\label{ch:sql-foundations}

\begin{learningobjectives}
After completing this chapter, you will be able to:
\begin{itemize}
  \item use \texttt{SELECT}, \texttt{FROM}, aliases, \texttt{LIMIT}, and
    \texttt{ORDER BY};
  \item filter with comparisons, logical operators, \texttt{IN},
    \texttt{BETWEEN}, \texttt{LIKE}, and \texttt{IS NULL};
  \item construct categories with \texttt{CASE};
  \item use \texttt{DISTINCT}, \texttt{COUNT}, \texttt{SUM}, \texttt{AVG},
    \texttt{MIN}, and \texttt{MAX};
  \item aggregate with \texttt{GROUP BY} and filter groups with
    \texttt{HAVING}; and
  \item state and validate the input and output grain of every query.
\end{itemize}
\end{learningobjectives}

\section{The shape of a query}

A SQL query is a request for a relation. Introductory queries often follow:

\begin{lstlisting}[language=SQL]
SELECT column_or_expression
FROM schema.table
WHERE row_condition
GROUP BY grouping_columns
HAVING group_condition
ORDER BY sorting_expression
LIMIT number_of_rows;
\end{lstlisting}

SQL is written in this conventional order, although the database reasons about
the clauses in a different logical sequence. Begin by identifying
\texttt{FROM}: it tells you the input table and therefore the input grain.

\section{Inspect a table safely}

\subsection{Question}
What columns and example values appear in the raw course table?

\begin{lstlisting}[language=SQL,caption={Instructional inspection query.}]
SELECT
    code_module,
    code_presentation,
    module_presentation_length
FROM raw.courses
ORDER BY code_module, code_presentation
LIMIT 5;
\end{lstlisting}

\subsection{Reasoning}
\begin{itemize}
  \item \textbf{Input grain:} one module-presentation.
  \item \textbf{Output grain:} one module-presentation, capped at five rows.
  \item \textbf{\texttt{SELECT}:} names the returned columns.
  \item \textbf{\texttt{FROM}:} names the schema-qualified table.
  \item \textbf{\texttt{ORDER BY}:} makes the limited sample deterministic.
  \item \textbf{\texttt{LIMIT}:} protects against accidentally printing a large
    table.
\end{itemize}

\begin{validationbox}
\texttt{LIMIT 5} does not prove that the table contains five rows or that the
sample represents the full distribution. Verify total row count separately.
\end{validationbox}

\section{Count rows and distinct entities}

\subsection{Question}
How many course rows, anonymized modules, attempts, and unique students exist?

\begin{lstlisting}[language=SQL]
SELECT
    count(*) AS module_presentation_rows,
    count(DISTINCT code_module) AS modules
FROM raw.courses;

SELECT
    count(*) AS student_module_attempts,
    count(DISTINCT id_student) AS unique_students,
    count(*) - count(DISTINCT id_student)
        AS additional_attempt_records
FROM core.student_attempts;
\end{lstlisting}

\verified{The expected results are 22 module-presentations, seven modules,
32,593 attempts, 28,785 unique students, and 3,808 additional attempt records.}

\subsection{Why the two counts differ}
\texttt{count(*)} counts rows. \texttt{count(DISTINCT id\_student)} counts
different nonnull identifiers. Their difference proves that some identifiers
appear in more than one attempt. It does not tell you whether the repeat is
another module, another presentation, or both.

\section{Filter rows}

\subsection{Comparisons and logical operators}

\begin{lstlisting}[language=SQL]
SELECT
    code_module,
    code_presentation,
    id_student,
    registration_day
FROM core.student_attempts
WHERE registration_day < 0
  AND code_module IN ('AAA', 'BBB')
ORDER BY registration_day
LIMIT 25;
\end{lstlisting}

The query asks for attempts registered before day zero in either module AAA or
BBB. \texttt{AND} requires both conditions. \texttt{IN} is a readable form of
multiple equality alternatives.

\subsection{\texttt{BETWEEN} and \texttt{LIKE}}

\begin{lstlisting}[language=SQL]
SELECT
    code_module,
    code_presentation,
    id_student,
    activity_day,
    click_count
FROM core.vle_interactions
WHERE activity_day BETWEEN 0 AND 6
  AND code_presentation LIKE '2014%'
ORDER BY code_module, code_presentation, id_student
LIMIT 100;
\end{lstlisting}

\texttt{BETWEEN 0 AND 6} includes both endpoints. The pattern
\texttt{'2014\%'} means text beginning with 2014. The query's input and output
grain remain interaction rows; filtering changes which rows appear, not what
one row represents.

\section{Work with null values}

\subsection{Question}
How many attempt records have no deprivation-band value?

\begin{lstlisting}[language=SQL]
SELECT count(*) AS missing_imd_band
FROM staging.student_info
WHERE imd_band IS NULL;
\end{lstlisting}

\verified{The saved quality profile reports 1,111 rows.}

Do not write \texttt{imd\_band = NULL}. SQL null represents unknown or absent
information and must be tested with \texttt{IS NULL} or
\texttt{IS NOT NULL}.

\section{Create an interpretable category with \texttt{CASE}}

\subsection{Question}
Can registration days be summarized into a few timing bands?

\begin{lstlisting}[language=SQL,caption={Instructional categorical
transformation.}]
SELECT
    CASE
        WHEN registration_day IS NULL THEN 'unknown'
        WHEN registration_day < 0 THEN 'before course start'
        WHEN registration_day BETWEEN 0 AND 13 THEN 'first two weeks'
        ELSE 'later'
    END AS registration_timing,
    count(*) AS attempts
FROM core.student_attempts
GROUP BY registration_timing
ORDER BY attempts DESC;
\end{lstlisting}

\texttt{CASE} is evaluated from top to bottom. The explicit null branch matters:
a null comparison is not true, so an omitted branch would place unknown values
in \texttt{ELSE}.

\begin{validationbox}
Sum the grouped \texttt{attempts}. It should equal the total attempt count,
32,593. If it does not, investigate whether rows were filtered or a group was
lost.
\end{validationbox}

\section{Aggregate outcomes}

\subsection{Question}
How many attempts ended in each recorded final result?

\begin{lstlisting}[language=SQL]
SELECT
    final_result,
    count(*) AS attempts,
    round(
        100.0 * count(*) / sum(count(*)) OVER (),
        1
    ) AS percent_of_attempts
FROM core.student_attempts
GROUP BY final_result
ORDER BY attempts DESC;
\end{lstlisting}

\begin{table}[htbp]
\centering
\caption{Verified recorded outcome counts}
\label{tab:outcome-counts}
\begin{tabular}{lrr}
\toprule
Final result & Attempts & Approximate percent \\
\midrule
Pass & 12,361 & 37.9 \\
Withdrawn & 10,156 & 31.2 \\
Fail & 7,052 & 21.6 \\
Distinction & 3,024 & 9.3 \\
\bottomrule
\end{tabular}
\end{table}

The input grain is one attempt. The output grain is one final-result category.
The grouped counts must sum to 32,593. Percentages may sum to 99.9 or 100.1
after rounding.

\section{Use numerical aggregates}

\subsection{Question}
What is the range and average of recorded nonnull assessment scores?

\begin{lstlisting}[language=SQL]
SELECT
    count(score) AS recorded_scores,
    min(score) AS minimum_score,
    max(score) AS maximum_score,
    round(avg(score), 2) AS mean_score
FROM core.assessment_submissions;
\end{lstlisting}

\begin{itemize}
  \item \texttt{count(score)} excludes null scores.
  \item \texttt{min} and \texttt{max} establish observed boundaries.
  \item \texttt{avg} excludes nulls rather than treating them as zero.
  \item The output grain is one summary row.
\end{itemize}

The project enforces recorded scores between 0 and 100. This query should not be
used to infer the number of expected submissions because absent submissions
have no row in this table.

\section{Group by module-presentation}

\subsection{Question}
Which module-presentations have at least 1,000 attempt records?

\begin{lstlisting}[language=SQL]
SELECT
    code_module,
    code_presentation,
    count(*) AS attempts
FROM core.student_attempts
GROUP BY code_module, code_presentation
HAVING count(*) >= 1000
ORDER BY attempts DESC;
\end{lstlisting}

\texttt{WHERE} filters input rows before grouping. \texttt{HAVING} filters
groups after aggregation. The output grain is one retained
module-presentation, not one attempt.

\section{Conditional aggregation}

\subsection{Question}
How many outcomes of each type occur within a module-presentation?

\begin{lstlisting}[language=SQL,caption={Selected production pattern from
\texttt{analytics.course\_outcome\_summary}.}]
SELECT
    code_module,
    code_presentation,
    count(*) AS student_attempts,
    count(*) FILTER (
        WHERE final_result = 'Pass'
    ) AS passes,
    count(*) FILTER (
        WHERE final_result = 'Withdrawn'
    ) AS withdrawals,
    round(
        avg((final_result = 'Withdrawn')::integer)::numeric,
        4
    ) AS withdrawal_rate
FROM core.student_attempts
GROUP BY code_module, code_presentation;
\end{lstlisting}

PostgreSQL's \texttt{FILTER} clause calculates multiple conditional aggregates
without changing the group grain. Casting the Boolean comparison to integer
turns false into 0 and true into 1; its average is therefore a proportion.

\section{Validation habits for introductory SQL}

For every query, record:

\begin{table}[htbp]
\centering
\caption{SQL reasoning template}
\label{tab:sql-template}
\begin{tabularx}{\textwidth}{L{0.27\textwidth}Y}
\toprule
Question & Required statement \\
\midrule
Analytical question & What decision or understanding motivates the query? \\
Input & Which table or view is read? \\
Input grain & What does one input row represent? \\
Filter & Which rows are included or excluded? \\
Output grain & What does one result row represent? \\
Null risk & Which expressions exclude or preserve nulls? \\
Duplication risk & Can one input row appear more than once? \\
Expected evidence & Which counts, ranges, or categories are known? \\
Validation & How will row counts or totals be reconciled? \\
Later use & Which analysis could consume the result? \\
\bottomrule
\end{tabularx}
\end{table}

\section{Knowledge check}

\begin{enumerate}
  \item Why should \texttt{ORDER BY} accompany a reproducible limited sample?
  \item What is the difference between \texttt{count(*)} and
    \texttt{count(score)}?
  \item When is \texttt{HAVING} used instead of \texttt{WHERE}?
  \item Why does filtering not automatically change table grain?
  \item What null error can occur in a \texttt{CASE} expression?
  \item How can grouped totals validate an aggregation?
\end{enumerate}

\begin{exercisebox}{Exercise 8.1: Inspect assessment definitions}
Write a query returning module, presentation, assessment type, due day, and
weight from \texttt{core.assessments}. Keep non-exam assessments with a known
due day between 0 and 90, order by module-presentation and due day, and limit
the display to 50 rows. State input grain, output grain, null rule, and expected
use.
\solutionref{sol:sql-inspection}
\end{exercisebox}

\begin{exercisebox}{Exercise 8.2: Validate outcome percentages}
Extend the outcome count query to return a percentage. Verify that counts sum
to 32,593 and explain why rounded percentages might not sum to exactly 100.
\solutionref{sol:sql-outcomes}
\end{exercisebox}

\begin{exercisebox}{Exercise 8.3: Profile missingness by module}
Write a query that returns one row per module with:
\begin{itemize}
  \item total attempts;
  \item attempts with missing \texttt{imd\_band}; and
  \item missing percentage.
\end{itemize}
Use conditional aggregation. State why this is descriptive and not evidence
that module membership caused missingness.
\solutionref{sol:sql-missing}
\end{exercisebox}

\begin{exercisebox}{Exercise 8.4: Find repeated student identifiers}
Write a query returning student IDs that occur in more than one attempt,
including their number of attempts. Order from most to least attempts and limit
the display to 25. Explain why this output must not be joined back to
interaction rows without a module-presentation key.
\solutionref{sol:sql-repeats}
\end{exercisebox}

\begin{commonmistakes}
\begin{itemize}
  \item Running \texttt{SELECT *} on the 10.6 million-row interaction table
    without a limit.
  \item Omitting \texttt{ORDER BY} and treating an arbitrary sample as stable.
  \item Using \texttt{= NULL}.
  \item Treating null scores as zero without an explicit analytical rule.
  \item Selecting nongrouped columns that do not define the output grain.
  \item Interpreting grouped associations as causes.
\end{itemize}
\end{commonmistakes}

\transferheading
The SQL clauses in this chapter transfer directly to most relational datasets.
What changes is the domain meaning: table grain, valid categories, missingness,
and validation totals. Carry the reasoning template with the syntax.

\begin{summarybox}
Introductory SQL begins with a question and an input grain. Selection,
filtering, categorization, distinct counting, numerical aggregation, grouping,
and group filtering produce relations at explicit output grains. Null behavior
and grouped reconciliation are part of the query, not afterthoughts. The next
SQL chapters will build joins and windows on this foundation.
\end{summarybox}
